\barrerapagina
\section{Resultados experimentales}\label{sec:resultados}
\subsection{Comparación común de entrenamiento, validación y prueba}
Las métricas de la tabla \ref{tab:metricas} proceden de los estados exportados, no de una ejecución nueva. La evaluación final contiene exactamente 1\,111 observaciones de enero a mayo de 2016: 558 negativas y 553 positivas. La alineación de identificadores, productos, meses y etiquetas fue comprobada contra los metadatos aprobados; las matrices y métricas se reproducen desde las predicciones existentes.
\begin{table}[htbp]
\centering\small
\caption{Métricas finales por partición (\%). Estados seleccionados en validación; no corresponden a la última época del perceptrón.}\label{tab:metricas}
\begin{tabular}{lrrrrrrr}\toprule
Método & Conjunto & Acc. & Prec. & Recall & Espec. & F1 & BA \\ \midrule
Perceptrón & Train & 64.623 & 66.815 & 62.972 & 66.396 & 64.837 & 64.684 \\
Perceptrón & Validación & 61.720 & 60.632 & 69.300 & 53.965 & 64.677 & 61.632 \\
Perceptrón & Test & 53.915 & 52.008 & 96.022 & 12.186 & 67.471 & 54.104 \\
Logística & Train & 65.840 & 66.544 & 68.472 & 63.013 & 67.494 & 65.742 \\
Logística & Validación & 61.149 & 62.812 & 56.810 & 65.589 & 59.660 & 61.199 \\
Logística & Test & 62.286 & 57.666 & 91.139 & 33.692 & 70.638 & 62.415 \\
Mayoritario & Train & 51.794 & 51.794 & 100.000 & 0.000 & 68.242 & 50.000 \\
Mayoritario & Validación & 50.571 & 50.571 & 100.000 & 0.000 & 67.172 & 50.000 \\
Mayoritario & Test & 49.775 & 49.775 & 100.000 & 0.000 & 66.466 & 50.000 \\
Histórico & Train & 49.272 & 51.540 & 34.425 & 65.223 & 41.279 & 49.824 \\
Histórico & Validación & 49.049 & 49.453 & 34.011 & 64.434 & 40.303 & 49.222 \\
Histórico & Test & 51.395 & 51.884 & 32.369 & 70.251 & 39.866 & 51.310 \\
\bottomrule\end{tabular}
\end{table}

\figura{metricas_globales.png}{Métricas globales de test sobre los mismos casos. Todas son proporciones o porcentajes según la escala rotulada; ninguna mide beneficio económico.}{fig:metricas}{1}
La logística tiene la mayor BA global observada: 62,415\%, 8,311 puntos sobre el perceptrón, 12,415 sobre mayoritario y 11,106 sobre histórico. El perceptrón supera esos baselines en 4,104 y 2,794 puntos de BA. Son diferencias descriptivas de la reserva final, sin prueba de superioridad universal ni significación estadística.

El baseline mayoritario alcanza recall 100\% sin reconocer un solo negativo. Esto demuestra por qué accuracy o sensibilidad aisladas no resumen el equilibrio entre clases. El histórico obtiene especificidad 70,251\%, pero recall 32,369\%, y BA cercana a 50\%.

\subsection{Matrices de confusión, sensibilidad y exceso de positivos}
\begin{table}[htbp]
\centering\small
\caption{Conteos de test y tasas de error (\%). Filas reales y columnas predichas en orden 0,1.}\label{tab:confusiones}
\begin{tabular}{lrrrrrrr}\toprule
Método & TN & FP & FN & TP & FPR & FNR & Pred. 1 \\ \midrule
Perceptrón & 68 & 490 & 22 & 531 & 87.814 & 3.978 & 1021 \\
Logística & 188 & 370 & 49 & 504 & 66.308 & 8.861 & 874 \\
Mayoritario & 0 & 558 & 0 & 553 & 100.000 & 0.000 & 1111 \\
Histórico & 392 & 166 & 374 & 179 & 29.749 & 67.631 & 345 \\
\bottomrule\end{tabular}
\end{table}

\figura{confusion_perceptron.png}{Matriz de confusión de test del perceptrón. Filas: etiqueta real; columnas: clase predicha. La elevada sensibilidad coexiste con 490 falsos positivos.}{fig:conf-p}{0.66}
\figura{confusion_logistica.png}{Matriz de confusión de test de logística con umbral congelado 0,5. Reconoce 188 negativos y 504 positivos, con 370 FP y 49 FN.}{fig:conf-l}{0.66}
\figura{sensibilidad_especificidad.png}{Sensibilidad y especificidad de los cuatro métodos en test. Logística mejora el reconocimiento negativo frente al perceptrón, pero conserva especificidad limitada.}{fig:sensibilidad}{0.95}
El perceptrón predice 1 en 1\,021 casos (91,899\%) y logística en 874 (78,668\%), frente a 553 positivos reales (49,775\%). Su exceso es de 42,124 y 28,893 puntos porcentuales, respectivamente. Predecir positivo frecuentemente captura muchos positivos reales, pero también eleva FP. La precisión positiva es solo 52,008\% en perceptrón y 57,666\% en logística. F1 positivo no garantiza reconocimiento de la clase negativa.

La logística reconoce 120 negativos adicionales frente al perceptrón, con 27 positivos reconocidos menos. Así disminuye 93 errores totales y mejora BA, calculada con denominadores específicos de cada clase. Esta explicación aritmética no establece una causa empresarial de la diferencia.

\subsection{Generalización temporal y error por partición}
La BA test menos validación cambia $-7,528$ puntos en perceptrón y $+1,216$ en logística. El aumento global de logística no acredita estabilidad en cada mes. Las proporciones positivas reales son 51,794\% en train, 50,571\% en validación y 49,775\% en test, mientras las predicciones positivas del perceptrón pasan de 48,815\% a 57,801\% y 91,899\%; las de logística, de 53,294\% a 45,738\% y 78,668\%. Hay un cambio descriptivo marcado en las decisiones, con prevalencia global relativamente próxima.
\figura{error_particiones.png}{Error de clasificación final, definido como $1-\mathrm{Accuracy}$, por método y partición. No es pérdida logística ni conteo de actualizaciones. No existe una curva de test por época.}{fig:error-particiones}{1}
\begin{table}[htbp]
\centering\small
\caption{Todos los meses de test. Tasas en \%; soporte negativo/positivo.}\label{tab:mensual}
\begin{tabular}{lrrrrrr}\toprule
Mes & 0/1 & BA perc. & BA log. & Prec. log. & Recall log. & Espec. log. \\ \midrule
2016-01 & 125/94 & 55.740 & 63.753 & 52.632 & 85.106 & 42.400 \\
2016-02 & 155/64 & 54.667 & 65.071 & 43.299 & 65.625 & 64.516 \\
2016-03 & 81/138 & 51.235 & 56.065 & 66.019 & 98.551 & 13.580 \\
2016-04 & 107/120 & 51.238 & 56.012 & 56.281 & 93.333 & 18.692 \\
2016-05 & 90/137 & 52.620 & 51.127 & 60.909 & 97.810 & 4.444 \\
\bottomrule\end{tabular}
\end{table}

\figura{ba_mensual.png}{BA de los cuatro métodos en todos los meses de test. No se excluyen meses difíciles ni se selecciona un subperíodo por rendimiento.}{fig:ba-mensual}{1}
En logística, la especificidad desciende de 64,516\% en febrero a 4,444\% en mayo, mientras el recall de mayo es 97,810\%. Su BA mensual cae hasta 51,127\%. El perceptrón supera a logística en BA de mayo (52,620\%), a pesar de la ventaja global de esta última. El desempeño cambia con el período y no se resume correctamente con una declaración universal de superioridad.

\subsubsection{Por qué BA global y promedio mensual difieren}
Si $P_t=TP_t+FN_t$ y $N_t=TN_t+FP_t$, al agrupar los conteos se obtiene
\begin{align}
TPR_G&=\sum_t\frac{P_t}{\sum_sP_s}TPR_t,&TNR_G&=\sum_t\frac{N_t}{\sum_sN_s}TNR_t,\\
BA_G&=\tfrac12(TPR_G+TNR_G),&BA_M&=\frac15\sum_t BA_t.
\end{align}
Por tanto,
\begin{equation}BA_G-BA_M=\frac12\sum_t\left[\left(\frac{P_t}{\sum_sP_s}-\frac15\right)TPR_t+\left(\frac{N_t}{\sum_sN_s}-\frac15\right)TNR_t\right].\end{equation}
Los pesos son distintos para positivos y negativos; no basta ponderar cada BA mensual por el total de filas. En logística, los soportes positivos mensuales son 94, 64, 138, 120 y 137; los negativos son 125, 155, 81, 107 y 90. La agrupación produce recall 0,911392405 y especificidad 0,336917563. De ellos resulta BA 62,415\%, frente a 58,406\% del promedio mensual sin ponderar. La diferencia de 4,009691 puntos se verifica con tolerancia $10^{-12}$ en puntos porcentuales, utilizando las contribuciones por mes registradas en el análisis temporal. No es una contradicción ni un error de cálculo.

\subsection{Productos conocidos y nuevos}
\begin{table}[htbp]
\centering\small
\caption{Comparación de cohortes en test, sin excluir los productos nuevos.}\label{tab:cohortes}
\begin{tabular}{lrrrrrrrr}\toprule
Cohorte & Método & N & 0/1 & TN & FP & FN & TP & BA (\%) \\ \midrule
conocidos & Perceptrón & 1095 & 552/543 & 62 & 490 & 14 & 529 & 54.327 \\
conocidos & Logística & 1095 & 552/543 & 184 & 368 & 46 & 497 & 62.431 \\
conocidos & Mayoritario & 1095 & 552/543 & 0 & 552 & 0 & 543 & 50.000 \\
conocidos & Histórico & 1095 & 552/543 & 389 & 163 & 367 & 176 & 51.442 \\
nuevos & Perceptrón & 16 & 6/10 & 6 & 0 & 8 & 2 & 60.000 \\
nuevos & Logística & 16 & 6/10 & 4 & 2 & 3 & 7 & 68.333 \\
nuevos & Mayoritario & 16 & 6/10 & 0 & 6 & 0 & 10 & 50.000 \\
nuevos & Histórico & 16 & 6/10 & 3 & 3 & 7 & 3 & 40.000 \\
\bottomrule\end{tabular}
\end{table}

Se consideran conocidos los nombres presentes en train, según la regla congelada. En test hay 219 conocidos con 1\,095 filas y ocho nuevos con 16 filas de abril y mayo. Se conservan todos en el global. Entre los nuevos, el perceptrón tiene precision 100\% sobre solo dos positivos predichos, pero omite ocho de diez positivos reales. Logística reconoce siete de diez positivos y cuatro de seis negativos. Los conteos son insuficientes para sostener estabilidad de porcentajes o una ventaja extrapolable a nuevos catálogos.

\subsection{Concordancias y desacuerdos}
Ambos modelos aciertan 554 casos y ambos fallan 374; solo perceptrón acierta 45 y solo logística 138. Las cuatro categorías son excluyentes y suman 1\,111. Coinciden en 928 decisiones y discrepan en 183: acuerdo no equivale a acierto. Los 93 aciertos exclusivos netos de logística explican la diferencia de accuracy. Las agrupaciones por mes y cohorte están disponibles en el archivo de acuerdos; no se usaron para construir un clasificador combinado.
\figura{acuerdos.png}{Descomposición de aciertos y errores compartidos y exclusivos por mes. Las categorías conservan los conteos de las mismas predicciones alineadas; los detalles por cohorte están en el archivo estructurado.}{fig:acuerdos}{1}
